\subsection{等价类；商集}

设 \(f\) 是一个函数， \(E\) 是其定义域， \(F\) 是其图。关系“ \(x \in E\) 且 \(y \in E\) 且 \(f(x) = f(y)\) ”是 \(E\) 上的一个等价关系，称为与 \(f\) 关联的等价关系。它等价于关系 \((\exists z)((x, z) \in F \text{ 且 } (y, z) \in F)\) ，也就是 \((\exists z)((x, z) \in F \text{ 且 } (z, y) \in \overline{F})\) ，因此其图为 \(\overline{F} \circ F\) .

我们现在将证明集合 E 上的每个等价关系 R 都具有这种形式。设 G 为 R 的图形。对于每个 \(x \in E\) （非空）集合 \(\mathrm{G}(x) \subset \mathrm{E}\) 称为x关于R的等价类；因此它是所有 \(y \in E\) 使得 \(R\{x, y\}\) 的集合。每个可以写成 \(\mathrm{G}(x)\) 对于某个 \(x \in E\) 的集合称为（关于R的）等价类。等价类中的一个元素称为该类的代表。关于R的等价类的集合（即所有形如 \(\mathrm{G}(x)\) 对于某个 \(x \in E\) 的对象的集合）称为E关于R的商集，记为E/R。定义域为E、目标为E/R的映射 \(x \to \mathrm{G}(x)\) ( \(x \in E\) 称为E到E/R的典范映射。

C55. 设R是集合E上的一个等价关系，并设p是E到E/R的典范映射。则

\[
\mathrm{R} \left\{x, y \right\} \Longleftrightarrow (p (x) = p (y)).
\]

根据上述记号，设x和y是E中的元素，使得

\[
(x, y) \in G.
\]

那么 \(x \in E\) 并且 \(y \in E\) 的图；让我们证明 \(G(x) = G(y)\) 。由于 \(y \in G(x)\) ，我们有（命题1） \(G(y) \subset (G \circ G)(x) = G(x)\) 另一方面，我们还有 \((y, x) \in G\) ，由此 \(G(x) \subset G(y)\) 因此

\[
\mathbf {G} (x) = \mathbf {G} (y),
\]

即， \(p(x) = p(y)\) 。反之，若 \(G(x) = G(y)\) ，我们有 \(y \in G(y) = G(x)\) ，因此 \((x, y) \in G\) 。这就完成了证明。

典范映射 \(p\) （从 \(\mathbf{E}\) 到上 \(\mathbf{E} / \mathbf{R}\) ）（§ 3，第 8 号，定义 11）的一个截面，更简称为 \(\mathbf{E}\) 的截面（关于关系 \(\mathbf{R}\) ).

\minorheading{示例}

\begin{enumerate}
\def\labelenumi{(\arabic{enumi})}
\item
  设 R 为集合 E 上的等价关系“ \(x \in E\) 并且 \(y \in E\) 并且 \(x = y\) ”。元素 \(x \in E\) 的等价类即为集合 \(\{x\}\) ，而典范映射 \(x \to \{x\}\) 将 E 映到 E/R 上是双射的。
\item
  设 E、F 为两个集合，并设 R 为上的等价关系， \(E \times F\) 与映射相关联， \(pr_{1}\) 的 \(E \times F\) 到 E 上。关于 R 的等价类为形如 \(\{x\} \times F\) ，其中 \(x \in E\) ；该映射 \(x \to \{x\} \times F\) 是 E 到 \((\mathbf{E} \times \mathbf{F}) / R\) .
\end{enumerate}

设 R 是集合 E 上的一个等价关系。商集 E/R 是 \(\mathfrak{P}(E)\) 的一个子集，且 E/R 的恒等映射是 E 的一个划分（ \(\S4\) ，第 7 号）；因为若 G 是 R 的图，我们有 \(x \in \mathrm{G}(x)\) 对于所有 \(x \in E\) ，且若两个等价类 \(\mathbf{G}(x)\) 并且 \(\mathbf{G}(y)\) 有一个公共元素 \(z\) ，那么 \(\mathbf{R}\{x, z\}\) 并且 \(\mathbf{R}\{y, z\}\) ，因此 \(\mathbf{G}(x) = \mathbf{G}(y)\) 。此外，关系

\[
(\exists \mathrm{X}) (\mathrm{X} \in \mathrm{E} / \mathrm{R} \text { 且 } x \in \mathrm{X} \text { 且 } y \in \mathrm{X})
\]

等价于 R\{x, y\}.

\(\mathbb{J}\) 反之，设 \((X_{i})_{i\in I}\) 是集合 E 的一个划分。立即可验证，关系 \((\exists i)(i\in I \text{ 且 } x\in X_{i} \text{ 且 } y\in Y_{i})\) 是 E 上的一个等价关系 R；关于 R 的等价类恰为该划分中的各集合 \(X_{i}\) ，且映射 \(i\to X_{i}\) 是从 I 到 E/R 的双射。E 的子集 S，如果对每个 \(i\in I\) ，集合 \(S\cap X_{i}\) 都恰含一个元素，就称为关于 R 的等价类的代表系（或横截）。这个名称也用于指从某集合 K 到 E 的任意单射，只要 K 在该单射下的像是关于 R 的等价类的代表系；例如 E 关于 R 的任一截面。

